\section{Need-Sufficiency Architecture}
\label{sec:need_sufficiency}

The revealed-sacrifice observation channel
(\citep[Revealed~Sacrifice]{lasser2026paper8a}) admits only
sacrifices satisfying S1 (free choice with non-degrading
alternatives) into its aggregate, and on those S1-admissible
events the channel reads sacrifice with positive polarity:
more sacrifice equals more revealed value.  S1 is the gate; the
present paper supplies the structural condition that makes one
of S1's failure modes diagnosable.  Without that condition,
events that are structurally involuntary on a need dimension
would, if mistakenly admitted, be misread with positive polarity
when their actual polarity is reversed.  An agent who pays \$500 for insulin is not revealing
that insulin is worth \$500 beyond baseline, they are revealing
that baseline functionality costs \$500 to maintain.  The sacrifice
magnitude reveals \emph{cost-of-access} (how expensive it is for
this agent to reach minimum viable functionality), not
\emph{preference strength} (how much the agent values the outcome
beyond baseline).  Larger sacrifice means a worse situation, not
higher revealed value.

The framework needs a structural requirement on poset
construction that distinguishes need-satisfaction from want-pursuit
and identifies the boundary where sacrifice polarity flips.

\subsection{Three motivating observations}

\paragraph{Needs are neither purely maximizing nor minimizing.}
Every need has \emph{both} maximizing dimensions (quality, quantity,
reliability) and minimizing dimensions (cost, distance, time). Needs are not a
single-polarity class. What \emph{is} true: the minimizing dimensions of needs
are where the polarity problem lives, because cost-of-access is involuntary.

\paragraph{Scalar sufficiency hides structural defects.}
A single ``water access: sufficient'' measurement might hide that the water is
contaminated.  The agent has water (passes the scalar threshold) but the way
the need is met actively degrades future $\volL$ through downstream effects on
health, cognition, and reproduction.  The framework reports green while the
system corrodes from the root.

\paragraph{Downstream-safe bundles.}
Three levels of sufficiency, each strictly stronger:
\begin{itemize}
\item \emph{Level~1: Scalar threshold (insufficient).}
  ``Water: yes/no.''  Hides quality, reliability, and cost defects.
\item \emph{Level~2: Bundle threshold (necessary but not
  sufficient).}
  All dimensions of the need must independently meet their thresholds
        simultaneously. This catches contamination (quality dimension fails),
        but misses subtler failures where the satisfaction pathway degrades
        downstream capabilities over time. For example, when the contamination
        thresholds are incorrectly tuned.
\item \emph{Level~3: Downstream-safe bundle.}
  Bundle thresholds are all met \emph{and} the need-satisfaction pathway does
        not produce net-negative effects on capabilities in the need's
        downstream capability cone.
\end{itemize}

\subsection{Formal definitions}

\begin{definition}[Multi-Dimensional Need Bundle]
\label{def:need_bundle}
A capability $N$ identified as a need at position $p_N$ in the
poset has:
\begin{itemize}
\item \textbf{Dimensions} $\{d_1, \ldots, d_m\}$ with sufficiency
  thresholds $\{s_1, \ldots, s_m\}$.
\item Each dimension has a \textbf{polarity}: maximizing (quality,
  quantity, reliability) or minimizing (cost, distance, time).
\item \textbf{Downstream cone}
  $\Down(p_N) = \{c : c \text{ depends on } N \text{ in the
  poset}\}$.
\item \textbf{Satisfaction pathway} $\pi_N$: the specific way the
  agent meets the need (which source, at what quality, at what
  cost).
\end{itemize}
For cost dimensions (minimizing polarity), the benchmark
inverts via a domain-supplied positive cost reference
$\mathrm{cost}_k^{\mathrm{ref}} > 0$: define
$\mathrm{actual}_k = \mathrm{cost}_k^{\mathrm{ref}} /
\mathrm{cost}_k$ and the sufficiency threshold
$s_k = \mathrm{cost}_k^{\mathrm{ref}} /
\mathrm{cost}_k^{\mathrm{max}}$ where
$\mathrm{cost}_k^{\mathrm{max}}$ is the maximum acceptable
cost.  Both quantities are dimensionless ratios sharing the same
cost-reference normalization, so $\mathrm{actual}_k$ and $s_k$
live on a common scale and benchmark maximization aligns with
cost minimization.  The inversion is defined on
$\mathrm{cost}_k > 0$; the zero-cost limit is read as ``no
cost-dimension constraint binds'' and treated as
benchmark-saturated, with the convention $\mathrm{actual}_k =
\mathrm{cost}_k^{\mathrm{ref}} / \mathrm{cost}_k$ capped at
some domain-set ceiling $M_k > s_k$ to avoid the $\mathrm{cost}_k
\to 0^+$ singularity.  The reference and ceiling are domain
parameters; the paper does not legislate specific values.
\end{definition}

\begin{definition}[Downstream-Safe Sufficiency]
\label{def:downstream_safe}
Let $\Pi_N^0 = \{\pi : d_k(\pi) = s_k \text{ for all } k\}$ be
the class of pathways satisfying each dimension of~$N$ at exactly
its sufficiency threshold (assumed non-empty by governance
realizability hygiene; in discrete-valued domains where exact
threshold-achievement may be infeasible, $\Pi_N^0$ should be
relaxed to $\{\pi : d_k(\pi) \geq s_k \text{ and within a
domain-set tolerance}\}$, with the framework's downstream
arguments preserved by treating the relaxed set as the
admissible reference base).  Let
$\Pi_N^{\mathrm{ref}} \subseteq \Pi_N^0$ be a non-empty
\emph{admissible reference set}, designated externally by
governance (Remark~\ref{rem:reference_pathway}): the
threshold-level pathways against which downstream non-degradation
is evaluated.  Need $N$ at position $p_N$ is
\emph{downstream-safe sufficient} under satisfaction pathway
$\pi_N$ iff:
\begin{enumerate}
\item[(a)] \textbf{Bundle thresholds met:}
  $d_k(\pi_N) \geq s_k$ for all dimensions~$k$ (with cost
  dimensions inverted per Definition~\ref{def:need_bundle}).
\item[(b)] \textbf{Downstream non-degradation against all
  admissible references:} for each downstream capability $c \in
  \Down(p_N)$, let $\mathrm{quality}_c : \Pi_N^0 \cup
  \{\pi_N\} \to \mathbb{R}$ be a domain-supplied scalar quality
  function on satisfaction pathways, providing a total order on
  pathways relative to capability~$c$ (e.g., for medical care
  domain, the standard-of-care quality function defined on the
  bundle-of-care pathways; for water access,
  the contamination-adjusted access function).  The framework
  reads $\mathrm{quality}_c$ as a governance input, not a
  computed quantity; the choice of quality function is part of
  the institutional adjudication that
  Definition~\ref{def:downstream_safe} delegates to governance.
  Under such a function, $\pi_N$ dominates \emph{every}
  $\pi^* \in \Pi_N^{\mathrm{ref}}$ on every downstream
  capability simultaneously:
  \[
    \mathrm{quality}_c\bigl(\text{with } N \text{ satisfied via }
    \pi_N\bigr)
    \;\geq\;
    \mathrm{quality}_c\bigl(\text{with } N \text{ satisfied via }
    \pi^*\bigr)
    \qquad \text{for all } c \in \Down(p_N)
    \text{ and all } \pi^* \in \Pi_N^{\mathrm{ref}}.
  \]
\end{enumerate}
The universal quantifier over $\Pi_N^{\mathrm{ref}}$ is
load-bearing: a merely existential condition (``dominate
\emph{some} threshold-level pathway'') would admit materially
degrading pathways that happen to beat a pathologically weak
reference.  Under the common singleton case
$\Pi_N^{\mathrm{ref}} = \{\pi_N^{\mathrm{ref}}\}$, the condition
reduces to dominance against one canonical baseline.  Each
element of $\Pi_N^{\mathrm{ref}}$ must be a realizable pathway,
not a coordinatewise supremum over $\Pi_N^0$ that no single
pathway achieves.
\end{definition}

\begin{remark}[Governance designates the admissible reference
set]
\label{rem:reference_pathway}
The admissible reference set $\Pi_N^{\mathrm{ref}}$ is a public
governance input, not a framework-computed object.  Institutions
typically adopt a specific reference pathway as the canonical
baseline (``standard of care'' in medicine, ``best available
technology'' in environmental policy); in those settings
$\Pi_N^{\mathrm{ref}}$ is a singleton.  Multi-reference settings
arise when different disciplines canonicalize differently and
$\pi_N$ must satisfy all of them simultaneously; the
downstream-safety condition then requires robust dominance over
every admissible baseline.  Governance may update
$\Pi_N^{\mathrm{ref}}$ over time as standards evolve; the
framework reads the current $\Pi_N^{\mathrm{ref}}$ at evaluation
time and does not attempt to re-derive it from first principles.
\end{remark}

\begin{remark}[Structural connection to exercise-pathway machinery]
\label{rem:exercise_pathway}
The need-satisfaction pathway $\pi_N$ is \emph{structurally} an
exercise pathway in the sense
of~\citep[Definition~2]{lasser2026horizon}: a sequence of
capability invocations producing an outcome (here, the
satisfaction of need~$N$).  This identification lets the framework
treat $\pi_N$ as a first-class object rather than an abstract
``way of satisfying the need.''

The downstream-safety condition
(Definition~\ref{def:downstream_safe}(b)) is, however, a distinct
mathematical object from the concentration-risk measure
of~\citep[Definition~3]{lasser2026horizon}.  Concentration risk
is a scalar worst-case contraction estimate over an entire
pathway; downstream-safety is a per-capability quality
comparison $\mathrm{quality}_c(\pi_N) \geq \mathrm{quality}_c(\pi^*)$
against every reference pathway $\pi^* \in \Pi_N^{\mathrm{ref}}$,
indexed by $c \in \Down(p_N)$.  Neither reduces to the other.
Evaluating $\mathrm{quality}_c(\pi_N)$ in practice requires
either (i) direct observation of $c$'s performance conditional on
the pathway, or (ii) structural inference from the poset's
cooperative structure.  A full formalisation of the evaluation
procedure is outside this paper's scope; within the paper,
Definition~\ref{def:downstream_safe}(b) is a structural condition
the framework \emph{enforces} (a pathway failing the condition is
rejected as a satisfaction pathway for $N$), not a quantity the
framework \emph{computes} directly.

\emph{Coherence with the framework's privacy stance.}  Option~(i)
above (direct observation of $c$'s performance conditional on the
pathway) is precisely the kind of agent-interior measurement that
the companion paper's privacy commitments rule out at scale on the
sacrifice channel.  Downstream-safety evaluation therefore lives
in one of two regimes: (a) \emph{public structural} evaluation,
where $\mathrm{quality}_c$ is determined from governance-supplied
reference-pathway specifications and the poset's structural
cooperative/subsumption topology---no agent-interior access
required, suitable for population-scale deployment; (b)
\emph{consented-perimeter} evaluation, where institutions with
explicit consent (the regime of the companion paper's
exercise-indicator machinery,
\citep[Revealed~Sacrifice, Appendix~A]{lasser2026paper8a}) directly observe
pathway performance on enumerable downstream capabilities.
Mixed deployments use~(a) for poset-wide screening and~(b) for
adjudication of contested pathways within consented sub-domains.
The architecture is coherent with the framework's privacy stance
under either regime: regime~(a) requires no privacy commitment
because no interior data is read, and regime~(b) operates inside
the perimeter where privacy commitments are explicitly waived.
Outside both regimes---ad hoc agent-interior measurement at
population scale---downstream-safety evaluation would be both
infeasible and inconsistent with the framework's stated
deployment constraints, and the architecture should not be
deployed in such a regime.

A policy consequence, not a theorem: because need-satisfaction
pathways are not optional (the agent must satisfy the need), the
framework's response to a pathway failing
Definition~\ref{def:downstream_safe}(b) differs from its response
to a risk claim under~\citep{lasser2026wamura}.  A failing risk
pathway is \emph{restricted} (the capability is blocked).  A
failing satisfaction pathway must be \emph{replaced} (the need
must still be met, but through a pathway satisfying
Definition~\ref{def:downstream_safe}).  This governance
distinction is upstream of the paper's formal machinery; it names
the correct operational response when the condition fails, not a
new computation the framework performs.
\end{remark}

\begin{definition}[Sufficiency Threshold as Polarity Boundary]
\label{def:polarity_boundary}
The sufficiency bundle is the structural boundary where
revealed-sacrifice polarity flips, conditional on the
choice-set hypothesis stated below:
\begin{itemize}
\item \textbf{Below sufficiency on any dimension, with no
  third non-degrading alternative:} The agent is
  investing to reach minimum viable functionality and has no
  alternative path that avoids degrading some need.  Sacrifice
  is involuntary (the agent must reach sufficiency to operate
  and has no other option).  Sacrifice magnitude reveals
  \textbf{cost-of-access}---how expensive it is for this agent
  to reach baseline on this dimension.  Larger sacrifice =
  worse situation.
\item \textbf{Above sufficiency on all dimensions:} The agent
  has reached functionality on every need dimension; the
  need-insufficiency failure mode of S1 is removed.  Sacrifice
  reveals \textbf{preference strength} \emph{provided} the
  trade is voluntary in the residual S1 sense (free of
  coercion, duress, information asymmetry, and the other
  material confounds of S1; see Remark~\ref{rem:s1_above}).
  Above-sufficiency status alone does not establish S1 in
  full---it removes one structural failure mode while leaving
  others to institutional attestation.  Larger sacrifice =
  higher revealed value, conditional on residual S1 holding.
\item \textbf{Below sufficiency on some dimension, but a
  third non-degrading alternative exists:} polarity is not
  determined by sufficiency status alone.  S1 can still hold
  via the third alternative; classification requires
  institutional review (Definition~\ref{def:s1_sufficiency}
  and Remark~\ref{rem:s1_above}).
\end{itemize}
The sufficiency threshold is where the need-insufficiency
failure mode of Assumption~S1
(\citep[Revealed~Sacrifice]{lasser2026paper8a}) becomes structurally
diagnosable: events with retention-degradation \emph{and} no
third non-degrading alternative fail S1 by
Definition~\ref{def:s1_sufficiency}; above sufficiency, this
specific failure mode does not apply (though S1 can still fail
via other mechanisms---see
Definition~\ref{def:s1_sufficiency} and
Remark~\ref{rem:s1_above}).
\end{definition}

\begin{definition}[Polarity-Correct Benchmark]
\label{def:polarity_benchmark}
For each dimension $d_k$ of a need capability:
\begin{itemize}
\item \textbf{Below sufficiency:}
  $\mathrm{benchmark}_k = \min(\mathrm{actual}_k / s_k,\, 1)$.
  Climbs from $0$ toward $1$ as the agent approaches sufficiency.
  Caps at $1$.  Maximizing this equals closing the sufficiency
  gap.
\item \textbf{Above sufficiency:}
  $\mathrm{benchmark}_k = 1 + \alpha_k \cdot
  \max\!\bigl((\mathrm{actual}_k - s_k)/s_k,\, 0\bigr)$,
  where $\alpha_k < 1$ discounts above-sufficiency gains relative
  to below-sufficiency gains \emph{at the margin}: the
  derivative of $\mathrm{benchmark}_k$ with respect to
  $\mathrm{actual}_k$ is $1/s_k$ below sufficiency (where the
  benchmark is $\min(\mathrm{actual}_k/s_k, 1)$) and
  $\alpha_k/s_k$ above sufficiency, so the per-unit ratio of
  marginal contribution is $\alpha_k:1$ in favour of closing
  the gap rather than enhancing beyond it.  This is a
  marginal-rate claim, not an absolute-magnitude claim: because
  the above-sufficiency term is unbounded in
  $\mathrm{actual}_k$, sufficiently large above-sufficiency
  gains can dominate the closed-gap benchmark in absolute
  terms.  Bounding total above-sufficiency contribution
  requires either a domain-specific cap on
  $\mathrm{actual}_k/s_k$ or replacing the linear ramp with a
  saturating function (e.g., $\alpha_k \cdot \min(\cdot, K)$);
  this paper leaves the cap as an architecture parameter.
\end{itemize}
\end{definition}

Definition~\ref{def:polarity_benchmark} operates at the
single-capability level: it determines the benchmark score for
each dimension of a given capability.  It does not, by itself,
model the structural fact that need-satisfaction is a precondition
for the exercise of downstream capabilities: an agent without
water does not merely lose one bit of water-access capability,
they lose access to the entire class of cooperative capabilities
that require a hydrated body.  To capture this, the framework
imposes a multiplicative gate on $\volL$-aggregation through the
downstream cone.

\begin{definition}[Downstream-Gated Capability Weighting]
\label{def:downstream_gate}
For each capability $c \in P$, let $\mathcal{N}(c) = \{N : c \in
\mathrm{Down}(p_N)\}$ be the set of needs whose downstream cones
contain $c$.  For each such $N$, let $s_N(c) \in [0, 1]$ be the
sufficiency level of $N$ in the context where $c$ is exercised,
with $s_N(c) = 1$ when $N$ is downstream-safe sufficient
(Definition~\ref{def:downstream_safe}) and $s_N(c) = 0$ when $N$
is fully unmet.  In the common case where sufficiency is uniform
across the downstream cone, we abbreviate $s_N(c) \equiv s_N$ and
treat $s_N$ as a scalar.

The \emph{effective weight} of $c$ is:
\[
  w^{\mathrm{eff}}(c) \;=\; w(c) \cdot
  \prod_{N \in \mathcal{N}(c)} g\bigl(s_N(c)\bigr),
\]
where $g : [0, 1] \to [0, 1]$ is a \emph{gate function} satisfying:
$g(1) = 1$, $g(0) = \alpha_0 \in [0, 1)$ (the gate's
\emph{floor}), $g$ monotone non-decreasing, and $g$
\emph{continuous on $(0, 1]$ with right-limit
$\lim_{s \to 0^+} g(s) = \alpha_0$}.  Continuity on $(0, 1]$
excludes step-function gates with jump discontinuities in the
interior; the right-limit condition at $0$ pins the floor.
Right-continuity alone (e.g., a step from $\alpha_0$ at $s_0$
defined-from-the-right) would be insufficient because
right-continuous step functions admit interior jumps; we use
the stronger continuous-on-$(0,1]$ condition to rule them out.
A gate with $\alpha_0 = 0$ is called a
\emph{zero-floor gate}; a gate with $\alpha_0 > 0$ is called a
\emph{positive-floor gate}.  The four candidate gates of
Remark~\ref{rem:gate_candidates} all satisfy
continuity-on-$(0,1]$ with the prescribed right-limit at zero;
gates with interior jump discontinuities (e.g., bare step
functions) are excluded.  Linear, ReLU, and power-law
gates are zero-floor; leaky-ReLU with parameter $\alpha > 0$ is
positive-floor with $\alpha_0 = \alpha$.

The effective capability volume is:
\[
  \volL^{\mathrm{eff}}(g) \;=\; \sum_{c \in P} w^{\mathrm{eff}}(c).
\]
The summation is over the same capability census~$P$ that
defines $\volL$ in~\citep[Proposition~1]{lasser2026poset};
$\volL^{\mathrm{eff}}$ is constructed as a downstream-gated
\emph{reweighting} of $\volL$'s atomic capability sum, with
each $c \in P$ contributing $w^{\mathrm{eff}}(c) = w(c) \cdot
\prod_{N \in \mathcal{N}(c)} g(s_N(c))$ in place of $w(c)$.
We treat $\volL^{\mathrm{eff}}$ as a \emph{gated diagnostic
quantity} for need-context optimization, not as a new
poset-measure inheriting M1--M6 by construction.  M5
non-negativity is preserved trivially (each gate factor is in
$[0,1]$, so $w^{\mathrm{eff}}(c) \geq 0$).  Compatibility with
M4-type additivity on poset-independent direct sums and with
M6-type cooperative superadditivity depends on whether the
gate factors $\prod_N g(s_N(c))$ are constant across the
relevant subset (so the gated weights factor cleanly); for
homogeneous-sufficiency settings (Section
\ref{sec:threshold_setting}) the factors are constant on each
downstream cone and the additivity carries through, but the
general case requires care and we do not claim
$\volL^{\mathrm{eff}}$ inherits M4 or M6 unconditionally.  The
framework optimizes $\volL^{\mathrm{eff}}$ rather than raw
$\volL$ once needs have been identified, treating the gated
quantity as an objective-substitute justified by its
incentive properties (Proposition~\ref{prop:near_total_collapse},
Proposition~\ref{prop:need_dominance}) rather than by axiom
inheritance.
\end{definition}

\begin{remark}[Candidate gate functions]
\label{rem:gate_candidates}
The choice of $g$ is a design parameter that trades off
smoothness, computational simplicity, and the strength of
need-dominance.  Four candidates:
\begin{itemize}
\item \textbf{Linear:} $g(s) = s$.  Simplest; gentle
  need-dominance; $s_N(c) = 0.5$ halves the downstream
  contribution.
\item \textbf{Leaky-ReLU:} $g(s) = \alpha$ for $s \leq s_0$,
  then linear ramp from $\alpha$ at $s_0$ to $1$ at $s = 1$.
  Parameters: threshold $s_0 \in [0, 1)$, leak $\alpha \in [0, 1)$.
  Below threshold, downstream capabilities retain a constant
  floor $\alpha$; above, rapid ramp to full contribution.
\item \textbf{ReLU:} Leaky-ReLU with $\alpha = 0$.  Below
  threshold, downstream capabilities are fully suppressed; above,
  linear ramp.  Hardest gate; strongest need-dominance.
\item \textbf{Power law:} $g(s) = s^\gamma$ for $\gamma \geq 1$.
  Smooth everywhere; gradient vanishes at $s = 0$ for $\gamma > 1$,
  creating an \emph{optimization dead zone near zero
  sufficiency}: an agent with severely unmet need would receive
  near-zero marginal incentive from $\volL^{\mathrm{eff}}$ to
  improve sufficiency, despite the situation being where the
  framework most wants improvement.  This is a structural
  trade-off of high-$\gamma$ power-law gates: severity is
  weighted strongly in absolute level (collapse magnitude is
  large) but weakly in marginal gradient (rate of improvement
  is small at low $s$); domains where this dead zone is
  unacceptable should use leaky-ReLU or low-$\gamma$ power-law
  gates.  Power law also gives progressively stronger penalty
  as the sufficiency gap
  widens.  Single tunable parameter.
\end{itemize}
The appropriate gate shape depends on the domain.  Biological
needs with cliff-like failure modes (dehydration, oxygen
deprivation) warrant ReLU or high-$\gamma$ power law.  Needs with
graceful degradation (social connection, mild nutritional
deficit) warrant leaky-ReLU or low-$\gamma$ power law.  The paper
leaves $g$ as a domain-specific design parameter and states
subsequent results for any valid $g$.
\end{remark}

\begin{proposition}[Near-Total Collapse Under Unmet Need]
\label{prop:near_total_collapse}
Let $g$ be any valid gate function
(Definition~\ref{def:downstream_gate}) with floor $\alpha_0 =
g(0) \in [0, 1)$, and suppose $s_N(c) \equiv s_N$ for all $c \in
\Down(p_N)$ (homogeneous sufficiency).  Partition the poset as
$P = \Down(p_N) \sqcup \overline{\Down(p_N)}$, with weight sums
$W_{\mathrm{down}} = \sum_{c \in \Down(p_N)} w(c)$ and
$W_{\mathrm{out}} = \sum_{c \in \overline{\Down(p_N)}} w(c)$, so
$\volL^{\mathrm{raw}} = W_{\mathrm{down}} + W_{\mathrm{out}}$.
If the agent is at sufficiency level $s_N < 1$:
\[
  \volL^{\mathrm{eff}}(g)
  \;\leq\; g(s_N) \cdot W_{\mathrm{down}} + W_{\mathrm{out}}.
\]

\medskip\noindent
\textbf{Collapse limits (upper-bound form).}
As $s_N \to 0$, $g(s_N) \to \alpha_0$ by the floor condition.
The proposition's upper bound becomes:
\begin{itemize}
\item \emph{Zero-floor gate} ($\alpha_0 = 0$): the upper bound
  on the downstream contribution vanishes,
  $\volL^{\mathrm{eff}} \leq W_{\mathrm{out}}$ in the limit.
  When $W_{\mathrm{down}} \gg W_{\mathrm{out}}$, the upper
  bound on effective volume falls to a small fraction of the
  unrestricted value (and the actual effective volume is at
  most this).
\item \emph{Positive-floor gate} ($\alpha_0 > 0$): the upper
  bound on downstream contribution falls to
  $\volL^{\mathrm{eff}} \leq \alpha_0 W_{\mathrm{down}} +
  W_{\mathrm{out}}$.  The effective volume retains \emph{at
  most} a fraction $\alpha_0$ of downstream weight even under
  fully unmet need; equality with $\alpha_0 W_{\mathrm{down}} +
  W_{\mathrm{out}}$ requires all other gates evaluating to $1$
  (otherwise the additional gate factors push the actual
  effective volume strictly below the upper bound).
\end{itemize}

\medskip\noindent
The multiplicative structure makes need-satisfaction a
precondition for $\volL$ to register the \emph{full} downstream
portion of the agent's capabilities; an agent cannot maximize
$\volL^{\mathrm{eff}}$ without closing sufficiency gaps on needs
whose downstream cones cover their capability-weight.  Whether
unmet need collapses the downstream contribution \emph{entirely}
or only to the floor $\alpha_0 W_{\mathrm{down}}$ is a design
choice governed by the gate's floor parameter.
\end{proposition}

\begin{proof}
Split the sum over $P$ by membership in $\Down(p_N)$.  For
$c \in \Down(p_N)$, the effective weight is
$w^{\mathrm{eff}}(c) = w(c) \cdot g(s_N) \cdot \prod_{N' \in
\mathcal{N}(c),\, N' \neq N} g(s_{N'}(c)) \leq w(c) \cdot g(s_N)$
since each additional gate factor is $\leq 1$.  For
$c \in \overline{\Down(p_N)}$, $N \notin \mathcal{N}(c)$, so
$w^{\mathrm{eff}}(c) = w(c) \cdot \prod_{N' \in \mathcal{N}(c)}
g(s_{N'}(c)) \leq w(c)$.  Summing the two cases:
\[
  \volL^{\mathrm{eff}}
  \leq g(s_N) \cdot \!\!\sum_{c \in \Down(p_N)}\!\! w(c)
       \;+\; \!\!\sum_{c \in \overline{\Down(p_N)}}\!\! w(c)
  = g(s_N) \cdot W_{\mathrm{down}} + W_{\mathrm{out}}.
\]
Right-continuity of $g$ at $0$ gives $\lim_{s_N \to 0^+} g(s_N)
= g(0) = \alpha_0$, so the downstream term tends to $\alpha_0
W_{\mathrm{down}}$: zero for zero-floor gates, the floor value
for positive-floor gates.
\end{proof}

\begin{remark}[The non-downstream floor and the gate floor]
\label{rem:nondownstream_floor}
Two floors appear at $s_N = 0$: the always-present
\emph{non-downstream} floor $W_{\mathrm{out}}$ (capabilities
not depending on $N$, which no gate can suppress), and the
gate-dependent \emph{downstream floor} $\alpha_0
W_{\mathrm{down}}$ (which is zero for zero-floor gates and
$\alpha \cdot W_{\mathrm{down}}$ for leaky-ReLU with leak
$\alpha > 0$).  The worked table in
Remark~\ref{rem:gate_comparison} exhibits both: at $s_N = 0$,
linear, ReLU, and power-law gates give $0.231 =
W_{\mathrm{out}}/\volL^{\mathrm{raw}} = 3/13$ (only the
non-downstream floor), while leaky-ReLU with $\alpha = 0.1$
gives $0.308 = (3 + 10 \cdot 0.1)/13 = (W_{\mathrm{out}} +
\alpha_0 W_{\mathrm{down}})/\volL^{\mathrm{raw}}$ (both floors).
The proposition's ``total collapse'' force applies to zero-floor
gates with $W_{\mathrm{down}} \gg W_{\mathrm{out}}$; for
positive-floor gates, the collapse is bounded below by the
gate's floor regardless of $W_{\mathrm{down}}/W_{\mathrm{out}}$
ratio.
\end{remark}

\begin{proposition}[Need-Dominance Marginal Rate, conditional
on a domain-supplied intervention metric]
\label{prop:need_dominance}
Assume homogeneous sufficiency $s_N(c) \equiv s_N$ for all $c \in
\Down(p_N)$ and that all other needs are at sufficiency
($s_{N'}(c) = 1$ for $N' \neq N$), so that $s_N$ is a scalar.
We compute partial derivatives of $\volL^{\mathrm{eff}}$
with respect to the design parameters $s_N$ and $w(c')$.
The dominance \emph{condition} below is conditional on a
domain-supplied intervention metric that prices unit
increments of $s_N$ and unit increments of $w(c')$ on a common
cost scale; without that metric, the partial derivatives are
not directly comparable and no formal dominance result
follows.  This is a design-convention statement, not a
free-standing inequality on the framework's primitives.
For any valid gate function $g$ differentiable at $s_N$ with
$g'(s_N) > 0$, the two relevant partial derivatives of
$\volL^{\mathrm{eff}}$ are
\[
  \frac{\partial \volL^{\mathrm{eff}}}{\partial s_N}
  \;=\; g'(s_N) \cdot W_{\mathrm{down}}(N),
  \qquad
  \frac{\partial \volL^{\mathrm{eff}}}{\partial w(c')}
  \;=\;
  \begin{cases}
    1 & c' \notin \Down(p_N), \\
    g(s_N) \leq 1 & c' \in \Down(p_N),
  \end{cases}
\]
where $W_{\mathrm{down}}(N) = \sum_{c \in \Down(p_N)} w(c)$.
The two derivatives have different domains ($s_N$ is a
dimensionless sufficiency level in $[0,1]$, while $w(c')$
carries $\volL$-weight units) and are not directly
commensurable.  A meaningful comparison requires a
domain-supplied \emph{intervention-cost map}
$\kappa : \Delta s_N \times \Delta w \to \mathbb{R}_{\geq 0}$
that prices increments in $s_N$ and increments in $w$ on a
common cost scale; the marginal-dominance discussion below is
the statement of when, under such a $\kappa$ identifying one
unit of $\Delta s_N$ with one unit of $\Delta w$ as equally
costly, the sufficiency gradient dominates.  Without~$\kappa$,
no comparison is meaningful; see Remark~\ref{rem:dominance_scale}.

\medskip\noindent
\textbf{Marginal-dominance condition (subject to a
domain-supplied intervention metric).}
``Closing the sufficiency gap is more valuable than adding
unit weight'' is a statement about a per-unit-of-intervention
comparison; it requires a domain-supplied metric that maps
sufficiency increments and weight increments to a common
intervention-cost scale.  Under such a metric, with one unit
of $\Delta s_N$ and one unit of $\Delta w$ identified as
equally costly, the marginal return to closing the sufficiency
gap on $N$ exceeds the marginal return to unit-weight addition
whenever
\[
  g'(s_N) \cdot W_{\mathrm{down}}(N) \;>\; 1.
\]
Without the metric, the partial derivatives are not
directly comparable and no marginal-dominance claim follows
from the formula alone.
For linear gates ($g(s) = s$), this reduces to
$W_{\mathrm{down}}(N) > 1$.  For leaky-ReLU or ReLU gates, the
condition holds at any $s_N$ in the ramp region provided
$W_{\mathrm{down}}(N)$ exceeds the reciprocal of the local
slope.  For power-law gates $g(s) = s^\gamma$ with $\gamma > 1$,
the condition fails near $s_N = 0$ (where $g'(s_N)$ is small)
even for large downstream cones---an intentional feature of
gates that weight severity of unmet need over its breadth.
\end{proposition}

\begin{proof}
Under the stated hypothesis, every gate product collapses: for
$c \in \Down(p_N)$, the product $\prod_{N' \in \mathcal{N}(c)}
g(s_{N'}(c))$ equals $g(s_N) \cdot \prod_{N' \neq N} g(1) =
g(s_N)$.  Thus
$w^{\mathrm{eff}}(c) = w(c) \cdot g(s_N)$ for
$c \in \Down(p_N)$ and $w^{\mathrm{eff}}(c) = w(c)$ otherwise.
Summing gives
$\volL^{\mathrm{eff}} = g(s_N) \cdot W_{\mathrm{down}}(N) +
W_{\mathrm{out}}(N)$.
Only the first term depends on $s_N$, so differentiation yields
$\partial \volL^{\mathrm{eff}}/\partial s_N =
g'(s_N) \cdot W_{\mathrm{down}}(N)$.  Differentiating in
$w(c')$: the capability $c'$'s gated weight is $g(s_N) \cdot
w(c')$ if $c' \in \Down(p_N)$ and $w(c')$ otherwise, giving the
stated partial derivatives.  Comparing the two partials on the
common unit-step scale yields the dominance condition.
\end{proof}

\begin{remark}[Finite-step form and intervention-cost scale]
\label{rem:dominance_scale}
The marginal-dominance condition is a per-unit-cost comparison
under the convention that one unit of $\Delta s_N$ and one unit
of $\Delta w$ have the same intervention cost.  Its finite-step
form is the Taylor expansion: a sufficiency increment
$\Delta s_N = \delta$ produces gain
$[g(s_N + \delta) - g(s_N)] \cdot W_{\mathrm{down}}(N) =
g'(s_N) W_{\mathrm{down}}(N) \cdot \delta + O(\delta^2)$, while
a weight increment $\Delta w(c') = \delta$ on a single
capability produces gain $\delta \cdot \max(1, g(s_N)) \leq
\delta$.  For small enough $\delta$, the marginal condition
$g'(s_N) W_{\mathrm{down}}(N) > 1$ implies the sufficiency
increment strictly dominates; a finite $\delta$ for which
$[g(s_N + \delta) - g(s_N)] \cdot W_{\mathrm{down}}(N) > \delta$
is the exact condition.  The condition
$g'(s_N) W_{\mathrm{down}}(N) > 1$ is \emph{not} a claim that a
finite $\Delta s_N$ of arbitrary size beats a finite unit of
weight: e.g., with $g(s) = s$ and $W_{\mathrm{down}} = 10$, a
$\Delta s_N = 0.01$ gain of $0.1$ does not exceed a
unit-weight gain of $1$ in absolute terms---only in
per-unit-of-intervention terms (the sufficiency increment
returned ten times the per-unit gain).  Interpreting the
condition as anything other than a marginal-rate comparison
would conflate the two intervention scales.
\end{remark}

\begin{remark}[Worked comparison across candidate gates]
\label{rem:gate_comparison}
Consider an agent with $10$ cooperative capabilities in
$\mathrm{Down}(p_N)$ (weight $1$ each) and $3$ capabilities
outside the cone (weight $1$ each), so
$\volL^{\mathrm{raw}} = 13$.  At homogeneous sufficiency level
$s_N(c) \equiv s$ for all $c \in \Down(p_N)$,
$\volL^{\mathrm{eff}}(s) = 3 + 10 \cdot g(s)$.  The following
table compares effective volume as a fraction of
$\volL^{\mathrm{raw}}$ for the four candidate gates (Leaky-ReLU
with $\alpha = 0.1$, $s_0 = 0.8$; ReLU with $s_0 = 0.8$; power
law with $\gamma = 2$):

\begin{center}
\small
\begin{tabular}{@{}cccccc@{}}
\toprule
$s$ & Linear & Leaky-ReLU & ReLU & Power ($\gamma{=}2$)
   & Power ($\gamma{=}3$) \\
\midrule
$0.00$ & $0.231$ & $0.308$ & $0.231$ & $0.231$ & $0.231$ \\
$0.25$ & $0.423$ & $0.308$ & $0.231$ & $0.279$ & $0.243$ \\
$0.50$ & $0.615$ & $0.308$ & $0.231$ & $0.423$ & $0.327$ \\
$0.75$ & $0.808$ & $0.308$ & $0.231$ & $0.663$ & $0.556$ \\
$0.80$ & $0.846$ & $0.308$ & $0.231$ & $0.723$ & $0.625$ \\
$0.90$ & $0.923$ & $0.654$ & $0.615$ & $0.854$ & $0.792$ \\
$1.00$ & $1.000$ & $1.000$ & $1.000$ & $1.000$ & $1.000$ \\
\bottomrule
\end{tabular}
\end{center}

Linear gives the gentlest gate (at $s = 0.5$, downstream
contributions are only halved, so $62\%$ of raw $\volL$ is still
registered).  ReLU gives the hardest gate (below threshold
$s_0 = 0.8$, downstream contributions are fully suppressed;
$\volL^{\mathrm{eff}}$ stays at the non-downstream floor of
$3/13 \approx 23\%$).  Power-law gates sit between, with
higher $\gamma$ approaching ReLU's sharpness.  Leaky-ReLU
preserves the ReLU shape above threshold while leaving a small
floor $\alpha$ below.
\end{remark}

\subsection{Sacrifice-polarity (S1) connection}

\begin{definition}[Below-sufficiency sacrifice]
\label{def:s1_sufficiency}
Let $\mathcal{A}_i$ denote agent $i$'s choice set at trade
time, including the trade ($Y_n$ for $X_n$), the retention
alternative (refuse the trade, keep $X_n$), and any other
alternative actions or trades available.  A sacrifice event
$(i, X_n, Y_n, t_n)$ is a \emph{below-sufficiency sacrifice}
when both:
\begin{enumerate}
\item[(H1)] retention is degrading: keeping $X_n$ leaves the
  agent below the sufficiency threshold $s_k$ on some need
  dimension $d_k$;
\item[(H2)] no alternative in $\mathcal{A}_i$ other than the
  trade itself is non-degrading on every need dimension; that
  is, for every $a \in \mathcal{A}_i \setminus \{(\text{the
  trade})\}$, there exists some need dimension $d_{\ell(a)}$
  with threshold $s_{\ell(a)}$ such that taking action $a$
  leaves the agent below $s_{\ell(a)}$ on $d_{\ell(a)}$ (the
  failing dimension may differ across alternatives, and need
  not coincide with the $d_k$ from H1).
\end{enumerate}
\end{definition}

\begin{remark}[Below-sufficiency sacrifices fail S1]
\label{rem:s1_sufficiency_failure}
A below-sufficiency sacrifice fails S1
(\citep[Revealed~Sacrifice]{lasser2026paper8a}) directly: H2 asserts that no
alternative-to-the-trade in $\mathcal{A}_i$ is non-degrading,
which is exactly the S1 failure condition (S1 requires the
existence of at least one non-degrading alternative to the
trade).  H1 in this construction names \emph{why} the
alternatives are degrading---retention is one specific
degrading option---ensuring the failure is need-driven rather
than artifactual.  The named-pattern value of
Definition~\ref{def:s1_sufficiency} is operational: it picks
out the structural sub-class of S1 failures that the
need-sufficiency architecture identifies (retention is
degrading and no third option escapes degradation), letting
the institutional attestation layer filter such events
without rederiving the S1 logic per case.
Below-sufficiency \emph{targeting} (sacrifice aimed at
improving an unmet need dimension where H1 holds but H2 fails
because $\mathcal{A}_i$ contains a non-degrading non-trade
alternative) does not satisfy
Definition~\ref{def:s1_sufficiency} and does not imply S1
failure: S1 can still hold via that third alternative, in
which case the trade reveals genuine preference between
non-degrading options.
\end{remark}

\noindent\textbf{Consequence for the sacrifice channel.}
Below-sufficiency sacrifices in the proposition's narrow sense
(retention degrades \emph{and} no third non-degrading
alternative exists) are filtered from the revealed-sacrifice
aggregate by S1 before they reach the aggregate lower bound
(the aggregate lower bound
theorem~\citep[Revealed~Sacrifice]{lasser2026paper8a}).  Below-sufficiency
\emph{targeting} where a third alternative exists is not
automatically filtered by this proposition: such events may
still satisfy S1 if the agent's choice set contains a
non-degrading option, and the institutional attestation layer
must determine which class an event belongs to.  For
above-sufficiency events (retention does not put the agent
below threshold), the need-insufficiency failure mode of S1
is removed structurally.  S1 may still fail for
above-sufficiency events via other mechanisms (coercion,
duress---see Remark~\ref{rem:s1_above} below).

\begin{remark}[S1 is not automatically satisfied above
sufficiency]
\label{rem:s1_above}
S1 can still fail for above-sufficiency events if the trade was
coerced by other mechanisms (economic duress, political pressure,
information asymmetry) even though the need-sufficiency condition
is met.  The sufficiency threshold resolves one structural failure
mode of S1 (the involuntary-need-satisfaction case); other S1
failure modes (economic duress, information asymmetry, and the
coercion open question discussed in the companion
paper~\citep[Revealed~Sacrifice]{lasser2026paper8a}) remain.
\end{remark}

\begin{remark}[Operational sourcing of the H1+H2 conditions]
\label{rem:s1_sufficiency_sourcing}
H1 (retention degrades) is checkable from the agent's
sufficiency state on each need dimension, which is computable
from the poset construction once thresholds are supplied.  H2
(no third non-degrading alternative) is the harder claim: it
quantifies over the agent's full choice set $\mathcal{A}_i$,
which is not directly observable from public trade data.  In
practice, H2 must be established by the institutional
attestation layer of the companion
paper~\citep[Revealed~Sacrifice]{lasser2026paper8a} (the external-attestation
remark there describes how per-event hypotheses about the
agent's choice set are certified through trade-level audit or
attestation, not through the public partition $\mathcal{C}$
alone).  Without H2 attestation, an event with H1 holding
should be classified as below-sufficiency \emph{targeting}
(potentially S1-admissible) rather than below-sufficiency
\emph{sacrifice} (S1-failing); the conservative default is to
require explicit H2 attestation before filtering.
\end{remark}

\subsection{The capability-stuffing problem}
\label{sec:capability_stuffing}

An agent optimizing $\volL$ under the non-triviality axiom M5
(every valid capability with $s_{\max} \geq 1$ contributes
positive weight~\citep[Proposition~1]{lasser2026poset}) has a
natural incentive to add capabilities that pass structural
validation but do not increase experiential optionality.  We call
this the \emph{capability-stuffing} failure mode.

The canonical form is redundant capability construction targeting
a single underlying need from multiple vantage points.  ``Travel
from home to a gas station'' is a real, benchmarkable
\emph{do}-capability: it satisfies the axioms M1--M6, has a
well-defined metric ($s_{\max} \geq 1$), and represents genuine
structural knowledge about the world.  Adding ``travel from work
to a gas station,'' ``travel from the grocery store to a gas
station,'' and so on produces a sequence of further real,
benchmarkable capabilities, each passing the same axiomatic
checks.  Collectively they inflate $\volL$ without broadening
what the agent can actually do: every addition instantiates the
same underlying optionality (fuel access) from a different
origin.

Capability-stuffing is distinct from other failure modes in the
sequence.  It is not \emph{wireheading}
(\citealp[Discussion]{lasser2026exo}), which is narrow
self-reinforcing exercise of existing capabilities; stuffing
adds new capabilities without exercising them degenerately.
It is not the \emph{Doll
Problem} (\citealp[Section~7.1]{lasser2026gfm}), which is
possessed-but-unexercised capabilities; the stuffed capabilities
may be technically exercised (routing through the gas station
``exercises'' the distance-measurement).  It is not the
\emph{B-to-C gap} in the aggregate sense this paper otherwise
treats; stuffing widens the gap through a specific mechanism
(individual-weight inflation) that the aggregate lower bound of
the aggregate lower bound theorem~\citep[Revealed~Sacrifice]{lasser2026paper8a} does not, by itself, address.
Stuffing is a framework-native perverse instantiation: the
optimizer discovers it by following the $\volL$-maximization
gradient under M5, not by departing from the framework's rules.

\paragraph{Tension with structural discovery.}
Capability-stuffing is closely related to a \emph{legitimate}
mechanism: the structural discovery value
of~\citep[Proposition~2]{lasser2026horizon}, which holds that
discovering new poset structure is value-positive.  An agent
adding a new distance-measurement capability to its repertoire is
making a structural discovery in that sense.  The formal
distinction between legitimate discovery and capability-stuffing
turns on whether the new capability creates new cooperative
structure:

\begin{itemize}
\item \textbf{Legitimate structural discovery} enriches the
  poset by creating cooperative capabilities (axiom M6
  superadditivity): the new capability interacts with existing
  capabilities to produce outputs previously unachievable.  The
  marginal contribution to $\volL$ includes cooperative terms
  beyond the individual weight.
\item \textbf{Capability-stuffing} adds individual capabilities
  whose marginal cooperative contribution is zero (or
  effectively so): the new capability's weight accrues through
  M5 alone, not through M6.  No new cooperative outputs become
  achievable.
\end{itemize}

Without a structural mechanism that distinguishes the two, the
framework rewards both equally, and a $\volL$-maximizer with
capability-discovery latitude will exploit the undifferentiated
reward to stuff.

\paragraph{What the need-sufficiency architecture does resolve.}
The downstream-safe bundle sufficiency condition
(Definition~\ref{def:downstream_safe}) and the polarity-correct
benchmark (Definition~\ref{def:polarity_benchmark}) together
handle one specific class of redundant construction: capabilities
whose claimed contribution lies \emph{below} the sufficiency
threshold on an already-satisfied need dimension.  For those, the
below-sufficiency benchmark saturates at $1$, so additional
instances do not inflate the benchmark.  This is a real structural
property of the architecture and it is worth stating narrowly:

\begin{proposition}[Below-Sufficiency Saturation (pointwise)]
\label{prop:below_suff_saturation}
For any need dimension $d_k$ of any single capability, the
below-sufficiency component of the polarity-correct benchmark
(Definition~\ref{def:polarity_benchmark}) saturates at $1$ once
$\mathrm{actual}_k \geq s_k$.  This is a \emph{pointwise} cap:
the benchmark value assigned to dimension $d_k$ of one
capability does not exceed $1$, regardless of how many times
that capability's $\mathrm{actual}_k$ is re-evaluated or
re-instantiated.
\end{proposition}

\begin{proof}
Below-sufficiency benchmark for dimension $d_k$ is
$\min(\mathrm{actual}_k / s_k, 1)$, which is bounded above by $1$
pointwise.
\end{proof}

\begin{remark}[The pointwise cap does not extend to an
aggregate dimension-level cap]
\label{rem:pointwise_not_aggregate}
Proposition~\ref{prop:below_suff_saturation} is a
\emph{per-capability, per-dimension} cap, not a dimension-level
aggregate cap.  Distinct capabilities $c_1, c_2, \ldots$ each
reporting $\mathrm{actual}_k \geq s_k$ on the same dimension
$d_k$ each score $1$ under the below-sufficiency component;
their M5-weighted sum in $\volL$ is not capped by this
proposition and accumulates with each new capability.  The
proposition therefore blocks only a narrow class of stuffing:
re-instantiating a \emph{single} capability to inflate its own
benchmark via repeated reporting.  The broader class of
capability-stuffing---creating distinct capabilities that each
satisfy sufficiency on the same dimension---is not structurally
blocked and is acknowledged as unresolved below.
\end{remark}

\paragraph{What the need-sufficiency architecture does not
resolve.}
Proposition~\ref{prop:below_suff_saturation} is the full extent
of the architecture's structural anti-stuffing contribution.  It
does \emph{not} cover the following:
\begin{enumerate}
\item[(i)] Above-sufficiency re-instantiation.  An attacker can
  construct capabilities with distinct
  $\mathrm{actual}_k$ values exceeding $s_k$ (e.g.,
  home-to-gas at 15~minutes and work-to-gas at 5~minutes).
  Each earns positive above-sufficiency credit
  $\alpha_k(\mathrm{actual}_k - s_k)/s_k$ under
  Definition~\ref{def:polarity_benchmark}, even though the
  underlying optionality (fuel access) is unchanged.
\item[(ii)] Capabilities whose claimed dimension is not inside
  any existing need's downstream cone, for which the gate
  function (Definition~\ref{def:downstream_gate}) does not
  apply and M5 weight accrues fully to $\volL$.
\item[(iii)] Arbitrary asymmetric decoration --- distinct
  cooperative descendants, minor weight perturbations, unique
  subsumption edges --- that preserves M5 credit for the stuffed
  capability regardless of whether the capability provides
  genuinely new optionality.
\end{enumerate}
Each of these defeats any attempt to block stuffing using
post-declaration poset properties or exercise patterns alone.
The honest statement of the need-sufficiency architecture's
contribution is sacrifice-polarity resolution (via
Definition~\ref{def:s1_sufficiency} above, which is
load-bearing for the sacrifice channel) together with the
narrow below-sufficiency cap of
Proposition~\ref{prop:below_suff_saturation}.  Full structural
resolution of M5 stuffing requires machinery this paper does not
develop.

\paragraph{Capability-stuffing as structural avoidance.}
The failure pattern across attempts is that every signature over
post-declaration data is adversarially manipulable.  The move
that changes the framing is recognizing stuffing as a specific
instance of \emph{structural avoidance}, the pathology flagged
as an open gap in~\citep[Discussion]{lasser2026poset}: when the
framework rewards capability declarations additively, an agent
optimizing $\volL$ has an incentive to \emph{decline to discover}
that capabilities merge, factor, or are otherwise reducible.
Stuffing is the concrete manifestation: the agent keeps
``travel from home to gas station'' and ``travel from work to
gas station'' distinct in $P$ rather than merging them into the
canonical form ``access to fuel from multiple origins'' that a
proper structural discovery would produce.  Each unmerged
capability inflates the M5 count; the merged capability has the
same optionality content.

The positive dual is already in the sequence:
\citep[Proposition~2]{lasser2026horizon} establishes that
structural discovery is value-positive under conditions on the
discovered structure's information value.  Making the symmetric
claim---that refusal to discover is value-negative---requires
mechanism, not entailment.  The direction is a \emph{proof-burden
taxonomy} with \emph{claim-and-challenge back-pressure}.

\paragraph{The proof-burden direction.}
For each avoidance route, a corresponding burden the agent must
discharge to retain $\volL$ credit.  Three types:
proof-of-non-mergeability (for the stuffing-specific case),
proof-of-provenance (blocking declarations without underlying
structural work), and proof-of-use (blocking declarations that
remain persistently dormant outside the residual class $\ResS$
of \S\ref{sec:residual_class}).  Proof-burdens operate as
preprocessing on the poset before M5 is applied; the M5 formula
is preserved, but the measure is evaluated on the canonicalized
poset rather than the presented one.  Canonicalization itself is
new axiomatic content at the preprocessing layer.

Claims are temporal and contestable.  An agent's
proof-of-non-mergeability stands until a counter-argument
demonstrates a merger mechanism; the discoverer of that
mechanism receives structural-discovery credit
(\citep[Proposition~2]{lasser2026horizon}); the avoider loses
the M5 credit whose proof has been voided.  This makes avoidance
economically risky rather than economically stable.  Analogous
architectures exist in scientific consensus, patent challenge,
and proof-of-stake slashing; adapting any of them requires
substantive protocol and calibration work.

\paragraph{Scope of the follow-up.}
The sequence provides substrate for parts of this construction:
Paper~5's commitment protocol for timestamped claims, the
companion paper's sacrifice channel~\citep[Revealed~Sacrifice]{lasser2026paper8a}
for partial use-evidence, Paper~6's controlled relaxation as an
adjudication pattern.  Substrate is not
implementation: formalizing the canonicalization operator, the
three proof-burden types with admissible evidence classes, the
claim-and-challenge protocol with challenge-market robustness
(admission filtering against spam, identity weighting against
Sybil flooding, independence requirements against cartel
suppression), and the equilibrium-dominance theorem requires a
dedicated follow-up paper.  This paper states the direction and
leaves the formal work there.  The open-question entry in
\S\ref{sec:discussion} describes the follow-up scope in more
detail.

\subsection{How sufficiency thresholds are set}
\label{sec:threshold_setting}

The framework treats sufficiency thresholds as
\emph{governance-supplied inputs}: the formal apparatus reads
$\{s_k\}$ at evaluation time and does not derive them from
first principles within the paper's measurement-theoretic
machinery (consistent with the limitations statement of
Section~\ref{sec:discussion}).  What is structurally rigorous
in the architecture is the \emph{downstream-safety condition}
(Definition~\ref{def:downstream_safe}) that any candidate
$\pi_N$ must satisfy at a given threshold, not the
threshold-setting procedure itself.  Four upstream mechanisms
can produce the governance input, not mutually exclusive:

\paragraph{(a) Poset-structural.}
Capabilities with high fan-out (many downstream dependents) have
their sufficiency level set by the minimum required for downstream
capabilities to function.  Computable from poset topology.

\paragraph{(b) Empirical / cross-sectional.}
Observe the population's actual capability levels and identify the
threshold below which agents demonstrably cease functioning on
downstream capabilities.

\paragraph{(c) Revealed-sacrifice-derived (heuristic).}
Sufficiency is the level at which sacrifice polarity flips.
Below sufficiency, sacrifice is involuntary and increasing; above
sufficiency, sacrifice becomes voluntary and choice-driven.  The
transition point in the sacrifice data identifies the threshold
empirically.  This is the most framework-native option but has a
bootstrap problem: identifying the threshold requires classifying
events, but classifying events requires the threshold.  This
mechanism is heuristic, not theorem-derived: the
revealed-sacrifice channel does not by itself certify a
threshold value, only suggest a transition region for governance
to ratify.  Resolution in practice: start with
structural/normative thresholds (mechanisms~(a) and~(d)),
refine using sacrifice data.

\paragraph{(d) Normative / external.}
A human authority or governance process declares sufficiency
levels.  This is the approach of~\citet{sen1999development}
and~\citet{nussbaum2011creating}.  The framework's formal
apparatus is agnostic to the source; it needs the thresholds as
input.

\medskip\noindent
\textbf{Operational stance.}  Initial thresholds are set by
(a) (poset-structural) and~(d) (normative), which are
non-bootstrapping.  Once enough sacrifice events are observed,
(c) (revealed-sacrifice-derived) refines the initial values
empirically, with the polarity-transition region in the trade
data calibrating the threshold against agents' actual
behavioural breakpoint.  Mechanism~(b) (cross-sectional
empirical) plays a similar refinement role on the
population-level functioning side.  No single mechanism is
``primary''; (a)+(d) initialize, (b)+(c) refine.

\medskip\noindent
\textbf{Descriptive vs.\ normative thresholds.}
Mechanisms~(a)--(c) produce \emph{descriptive} thresholds:
they identify levels at which downstream functionality
breaks down or sacrifice polarity inverts---empirical or
structural facts about agents in the deployment context.
Mechanism~(d) produces a \emph{normative} threshold:
the level a human authority has \emph{decided} agents should
not have to fall below, irrespective of the descriptive
breakdown level.  Descriptive and normative thresholds can
diverge in either direction: a society may set a normative
food-sufficiency threshold above the biological minimum to
respect dignity, or (perversely) a normative threshold below
the descriptive breakdown level for budgetary reasons.  The
framework's formal apparatus reads whatever threshold
governance supplies and does not by itself adjudicate between
descriptive and normative provenance; downstream-safety and
the structural arguments of
Section~\ref{sec:need_sufficiency}
proceed identically under either kind of threshold.  However,
the \emph{interpretation} of $\delta_{\mathrm{covered}}$ alarms
and the meaning of below-sufficiency event filtering depend on
which kind of threshold is in force: a poset-structural
threshold means ``below this level, downstream capabilities
literally cannot function''; a normative threshold means
``below this level, the governance regime has stipulated
agents should not have to sacrifice.''  Deployments should
document the provenance of each threshold so that diagnostic
output is read against the correct interpretive backdrop.
Mixed deployments---structural thresholds for biological needs
and normative thresholds for socially-constructed needs---are
operationally common, and the framework accommodates them as
long as each threshold's provenance is recorded with the
threshold value.
